\documentclass[12pt,a4paper,twoside]{report}
\usepackage{graphicx}
\usepackage[utf8]{inputenc}
\usepackage[portuguese]{babel}
\usepackage[margin=2.5cm]{geometry}
\usepackage[nottoc,numbib]{tocbibind}
\usepackage{pdfpages}
\usepackage{parskip}
\usepackage{url}
\usepackage{amsmath, amssymb, amsthm}
\usepackage{booktabs}
\usepackage[colorinlistoftodos]{todonotes}
\usepackage[acronym]{glossaries}
\makeglossaries

\usepackage{listings}
\renewcommand{\lstlistingname}{Listagem}
\renewcommand*{\lstlistlistingname}{Lista de Listagens}
\lstset{
 xleftmargin = 1cm,
 basicstyle=\footnotesize, 
 numbers=left, 
 captionpos=b,
 columns=fullflexible,
 breaklines=true, 
 frame=single
}

\begin{document}

\pagenumbering{roman}
\newacronym{pme}{PME}{Pequenas e Médias Empresas}
\newacronym{sabi}{SABI}{Sistema de Análise de Balanços Ibéricos}
\newacronym{rsf}{RSF}{Random Survival Forests}
\newacronym{coxph}{Cox PH}{Cox Proportional Hazards}
\newacronym{pwp}{PWP}{Prentice-Williams-Peterson}
\newacronym{dl}{DL}{Deep Learning}
\newacronym{ml}{ML}{Machine Learning}
\newacronym{rfd}{RFD}{Recurring Financial Distress}
\newacronym{cif}{CIF}{Cumulative Incidence Function}
\newacronym{ibs}{IBS}{Integrated Brier Score}
\newacronym{lstm}{LSTM}{Long Short-Term Memory}


\chapter*{Draft de Planeamento Estratégico (Nota 20)}
\textit{Este documento segue a estrutura oficial da tese para facilitar a transição de conteúdos.}

\tableofcontents 
\listoffigures 
\listoftables 

\printglossary[type=\acronymtype,title=Acrónimos,nonumberlist]

\chapter{Introdução}
\pagenumbering{arabic}
\section{Introdução e Pivot Estratégico}

Esta investigação transcende a visão binária e estática da falência empresarial, adotando uma perspetiva dinâmica e recorrente. O fenómeno central é o \textit{Recurring Financial Distress}, analisado através do ciclo \textbf{saúde-doença-recaída}.

\subsection{Diferenciação Crítica: Estados Transientes vs. Absorventes}
A literatura tradicional colapsa frequentemente diferentes tipos de eventos numa única variável de "distress". Para atingir a excelência preditiva, distinguimos três estados fundamentais:
\begin{itemize}
    \item \textbf{Estado 0 (Saudável):} Operação normal sem sinais de alerta.
    \item \textbf{Estado 1 (Distress Transiente):} Verificação de critérios operacionais recuperáveis (ex: Capital Próprio negativo ou EBITDA insuficiente para cobrir juros).
    \item \textbf{Estado 2 (Falência/Terminal):} Verificação de critérios de dissolução legal ou liquidação (Estado Absorvente).
\end{itemize}

\section{Hipóteses de Investigação}

\begin{description}
    \item[H1 (Feature Value):] A inclusão de variáveis de \textit{Corporate Governance} e Macroeconomia reduz significativamente o erro de calibração (Brier Score) e melhora a discriminação (C-Index) em comparação com modelos puramente financeiros baseados em rácios do SABI.
    \item[H2 (Superioridade Arquitetural):] Modelos de Deep Learning Sequencial (\textbf{DynamicDeepHit}) superam modelos estáticos (Cox, RSF) ao capturarem a memória financeira das PME através de trajetórias longitudinais.
    \item[H3 (Dinâmica da Recorrência):] Os determinantes da recaída (transição $0 \to 1$ após um episódio prévio) são estruturalmente diferentes dos determinantes do primeiro episódio de stress, justificando o uso de modelos estratificados (\textbf{Cox PWP}).
\end{description}

\section{Inovações Técnicas Centrais}
\begin{enumerate}
    \item \textbf{DynamicDeepHit:} Arquitetura com RNN/LSTM para processar sequências temporais e riscos concorrentes.
    \item \textbf{Cox PWP (Prentice-Williams-Peterson):} Modelação da recorrência com estratificação por episódio.
    \item \textbf{Calibração de Precisão:} Foco no \textit{Integrated Brier Score} (IBS) para medir a qualidade das probabilidades estimadas.
\end{enumerate}


\chapter{Estado da Arte e Trabalho Relacionado}
% Aqui entrarão as bases de sobrevivência e Deep Learning
\section{Modelação de Intensidade e Processos de Contagem}

Para modelar eventos recorrentes, abandonamos a função de sobrevivência simples $S(t)$ em favor da \textbf{Função de Intensidade} $\lambda(t | \mathcal{F}_{t^-})$, onde $\mathcal{F}_{t^-}$ representa a história da empresa até ao instante imediatamente anterior a $t$.

\subsection{Modelo de Prentice-Williams-Peterson (PWP-GT)}
Ao contrário do modelo de Andersen-Gill, o PWP permite que o risco mude com o número de episódios ($k$), sendo o modelo de eleição para testar a nossa hipótese H3:

\begin{equation}
    \lambda_k(t | Z) = \lambda_{0k}(t) \exp(\beta_k' Z)
\end{equation}

Onde:
\begin{itemize}
    \item $\lambda_{0k}(t)$ é o hazard base para o $k$-ésimo episódio de stress.
    \item $\beta_k$ representa os coeficientes específicos para o episódio $k$, permitindo verificar se a importância da \textit{Governance} ou \textit{Macro} se altera após o primeiro colapso.
\end{itemize}

\section{Riscos Concorrentes (Competing Risks)}

No contexto de PME, o evento terminal (Falência) atua como um risco concorrente para a recuperação. No modelo \textit{Multi-State}, a transição do Estado 1 (Distress) tem dois caminhos mutuamente exclusivos:
\begin{enumerate}
    \item \textbf{Transição $1 \to 0$:} Recuperação para a saúde financeira.
    \item \textbf{Transição $1 \to 2$:} Progressão para a falência total (Estado Absorvente).
\end{enumerate}

A função estatística correta para analisar este fenómeno é a \textbf{Função de Incidência Cumulativa (CIF)}:
\begin{equation}
    CIF_k(t) = P(T \le t, \text{Cause} = k)
\end{equation}
Esta função é superior ao estimador de Kaplan-Meier nestes cenários, pois não assume independência entre as causas de falha.

\section{Arquitetura Dinâmica do DeepHit}
\todo[inline,color=red!40]{PENDENTE: Descarregar Lee (2019) Dynamic-DeepHit para /articles}

O \textbf{DynamicDeepHit} estima a probabilidade discreta de falha para cada causa $k$ e cada intervalo de tempo $t$, utilizando a trajetória histórica $X_{1:T}$:
\begin{equation}
    y = f_{\theta}(X_{1:T})
\end{equation}
A rede neuronal é composta por uma \textbf{RNN/LSTM partilhada} que captura a dinâmica temporal, seguida de \textbf{MLPs causa-específicas} que estimam o risco para Distress vs. Falência.

% Add citations to your bibliography entries using \cite{key}
\cite{*}

\chapter{Desenho e Especificação}
\section{Estrutura de Intervalos (Counting Process)}

Para modelar a recorrência e os riscos concorrentes, o dataset deve ser reestruturado para o formato de \textbf{Processo de Contagem (t\_start, t\_stop)}. Cada empresa é representada por múltiplos intervalos que indicam a permanência ou transição entre estados.

\begin{table}[h]
\centering
\caption{Exemplo de Estrutura de Dados Multi-State}
\begin{tabular}{@{}ccccccc@{}}
\toprule
\textbf{NIF} & \textbf{t\_start} & \textbf{t\_stop} & \textbf{From\_State} & \textbf{To\_State} & \textbf{Status} & \textbf{Episode (k)} \\ \midrule
123 & 0 & 3 & 0 & 1 & 1 & 1 \\
123 & 3 & 5 & 1 & 0 & 1 & 1 \\
123 & 5 & 8 & 0 & 2 & 1 & 2 \\ \bottomrule
\end{tabular}
\end{table}

\section{Algoritmo de Definição de Estados}

A lógica de transição segue os critérios de negócio definidos:
\begin{enumerate}
    \item \textbf{Eventos Terminais:} Se o \textit{Status} legal indica falência (C3), transição imediata para o Estado 2 (Absorvente).
    \item \textbf{Distress Transiente:} Se ocorrer insolvência técnica (C1: Equity < 0) ou incapacidade de cobertura (C2: EBITDA < Juros), transição para o Estado 1.
    \item \textbf{Recuperação:} Se uma empresa em Estado 1 deixa de verificar C1 e C2, regressa ao Estado 0.
\end{enumerate}

\section{Prevenção de Data Leakage: Lagging e Censura}

\begin{itemize}
    \item \textbf{Lagging Temporal:} Todas as variáveis independentes ($X$) devem ser defasadas em 1 ano ($t-1$) em relação ao estado alvo ($y_t$), refletindo a disponibilidade real da informação contabilística.
    \item \textbf{Gestão de Censura:} Empresas que deixam de reportar dados sem indicação de falência são marcadas como censuradas à direita no seu último registo ativo.
\end{itemize}

\section{Paradigma Estatístico: Cox PWP-GT}

O modelo de Prentice-Williams-Peterson será a nossa âncora estatística para a recorrência.
\begin{itemize}
    \item \textbf{Estratificação:} O modelo será estratificado pelo número de episódios (\textit{strata}), permitindo que o hazard base varie em cada recidiva.
    \item \textbf{Variáveis Dependentes do Tempo:} Incorporação de indicadores macro dinâmicos (PIB, Inflação, Juros) como covariáveis que evoluem em cada intervalo.
\end{itemize}

\section{Ensemble Learning: Random Survival Forests (RSF)}

O RSF é utilizado para capturar as complexas interações não-lineares entre rácios financeiros e indicadores de \textit{Governance}.
\begin{itemize}
    \item \textbf{Splitting Rule:} Utilização do \textit{Log-rank} para maximizar a diferença estatística entre as curvas de sobrevivência dos nós filhos.
    \item \textbf{Interpretabilidade:} Aplicação de valores SHAP (\textit{SHapley Additive exPlanations}) para decompor a contribuição de cada rácio para o risco individual da PME.
\end{itemize}

\section{Deep Learning: DynamicDeepHit}

Este modelo representa a inovação central do projeto, ao introduzir a memória sequencial na análise de sobrevivência.
\begin{itemize}
    \item \textbf{Encoder LSTM:} Processa a sequência temporal de rácios financeiros dos últimos 5 anos, extraindo um \textit{embedding} que representa a trajetória de saúde da empresa.
    \item \textbf{Decoders Causa-Específicos:} Redes neuronais independentes para prever a probabilidade de transição para Distress (Recorrência) vs. Falência (Risco Concorrente).
    \item \textbf{Função de Perda Híbrida:} Combinação da \textit{Log-Likelihood} para o tempo do evento com uma \textit{Ranking Loss} para garantir a ordenação correta do risco.
\end{itemize}

\section{Features de Alta Ordem e Justificação Teórica}

Para além dos rácios financeiros tradicionais, a nossa arquitetura de dados inclui variáveis que capturam dimensões não lineares do risco.

\subsection{Distance to Default (Merton-inspired)}
\todo[inline,color=orange!40]{PENDENTE: Confirmar se Hillegeist (2004) está na pasta /articles}
Adaptamos a teoria de opções para PMEs sem cotação bolsista, utilizando a volatilidade do ROA como proxy da volatilidade dos ativos:
\begin{equation}
    DTD_i = \frac{\ln(A_i / D_i) + (\mu_A - \frac{\sigma_A^2}{2}) \cdot T}{\sigma_A \cdot \sqrt{T}}
\end{equation}

\subsection{Governance Entropy Index}
Mede a instabilidade da gestão como um sinal de alerta precoce (\textit{Early Warning}):
\begin{equation}
    GE_i(t) = -\sum_{k} p_k \ln(p_k)
\end{equation}
Onde $p_k$ é a proporção de tempo que o conselho de gerência manteve a mesma composição.

\subsection{Macro-Beta (Sensibilidade Heterogénea)}
Calcula a sensibilidade específica da PME ao ciclo económico nacional:
\begin{equation}
    \beta^{PIB}_{i} = \frac{Cov(\Delta Revenue_{i,t},\ \Delta PIB_t)}{Var(\Delta PIB_t)}
\end{equation}

\section{Estratificação de Feature Sets (H1)}

Para testar o valor incremental da \textit{Governance} e \textit{Macro} (H1), definimos três conjuntos de treino:
\begin{enumerate}
    \item \textbf{FS1:} Apenas rácios financeiros brutos (Baseline).
    \item \textbf{FS2:} Rácios financeiros + Indicadores de \textit{Governance}.
    \item \textbf{FS3:} Conjunto completo (Financeiros, Governance, Macro e Features de Alta Ordem).
\end{enumerate}


\chapter{Validação e Resultados (Planeamento)}
\section{Métricas de Discriminação}

Para avaliar a capacidade do modelo em ordenar corretamente as empresas por nível de risco:
\begin{itemize}
    \item \textbf{Harrell's C-Index:} Medida global de concordância entre os scores previstos e os tempos de falha observados.
    \item \textbf{Time-Dependent AUC:} Avaliação da performance discriminativa em horizontes temporais específicos (1, 3 e 5 anos), permitindo analisar a degradação da precisão ao longo do tempo.
\end{itemize}

\section{Métricas de Calibração (O Padrão de Excelência)}

Um modelo de nota 20 deve ser bem calibrado, não apenas discriminativo:
\begin{itemize}
    \item \textbf{Integrated Brier Score (IBS):} Métrica que quantifica a precisão das probabilidades estimadas ao longo de todo o período de estudo.
    \item \textbf{D-Calibration:} Teste estatístico para verificar se o modelo é consistente com a distribuição real dos dados de sobrevivência.
\end{itemize}

\section{Avaliação da Recorrência e Estabilidade}

\begin{itemize}
    \item \textbf{Mean Cumulative Function (MCF):} Comparação visual entre o número esperado de episódios de stress previsto vs. observado.
    \item \textbf{Stratified Evaluation:} Cálculo de métricas separadas para o primeiro episódio e para a recaída, testando a validade da hipótese H3.
\end{itemize}

\section{Protocolo de Validação Temporal}

\begin{itemize}
    \item \textbf{Train Set (2005-2018):} Base de aprendizagem histórica, incluindo o período da crise financeira.
    \item \textbf{Validation Set (2019-2021):} Tuning de hiperparâmetros, incluindo o choque do COVID-19.
    \item \textbf{Test Set (2022-2023):} Avaliação final e "cega" da performance preditiva em ambiente pós-pandémico.
\end{itemize}


\chapter{Conclusão}
% Resumo das metas de investigação

\bibliographystyle{ieeetr}
\bibliography{refs}

\end{document}
